% Mathematical target excerpt; Kauers and Koutschan (2023).
% Source: https://cs.uwaterloo.ca/journals/JIS/VOL26/Koutschan/kout4.tex
% Retained source notation; not a solution.
\begin{conjecture}\label{conj:kaprekar} If $(a_n)$ denotes the sequence \seqnum{A164735}, then
  for all $n\geq3$
\[
  a_{18k+i} = \frac{1}{40}\!
  \begin{cases}
    3 (243 k^5+405 k^4+35 k^3+395 k^2-318 k+40), & i=0; \\
    k (729 k^4-405 k^3-615 k^2+225 k+106), & i=1; \\
    729 k^5+1620 k^4+735 k^3+1320 k^2-684 k+40, & i=2; \\
    k (729 k^4-705 k^2+136), & i=3; \\
    3 k (243 k^4+675 k^3+515 k^2+565 k-118), & i=4; \\
    k (729 k^4+405 k^3-615 k^2-225 k+106), & i=5; \\
    3 k (243 k^4+810 k^3+845 k^2+790 k+32), & i=6; \\
    3 k (k+1) (243 k^3+27 k^2-142 k+12), & i=7; \\
    729 k^5+2835 k^4+3705 k^3+3405 k^2+726 k+40, & i=8; \\
    3 k (k+1) (243 k^3+162 k^2-127 k-18), & i=9; \\
    729 k^5+3240 k^4+5055 k^3+4860 k^2+1636 k+160, & i=10; \\
    3 k (k+1) (243 k^3+297 k^2-52 k-48), & i=11; \\
    729 k^5+3645 k^4+6585 k^3+6795 k^2+2926 k+400, & i=12; \\
    3 k (k+1) (243 k^3+432 k^2+83 k-58), & i=13; \\
    729 k^5+4050 k^4+8295 k^3+9270 k^2+4696 k+800, & i=14; \\
    3 k (k+1) (243 k^3+567 k^2+278 k-28), & i=15; \\
    3 (k+3) (243 k^4+756 k^3+1127 k^2+734 k+160), & i=16; \\
    3 k (k+1) (243 k^3+702 k^2+533 k+62), & i=17.
  \end{cases}
  \]
\end{conjecture}
